% Generated draft: requires mathematical and source review.
\documentclass[11pt]{article}
\usepackage[margin=25mm]{geometry}
\usepackage{fontspec}
\setmainfont{Latin Modern Roman}
\usepackage{amsmath,amssymb,mathtools,unicode-math}
\setmathfont{Latin Modern Math}
\usepackage{enumitem,needspace,xurl,hyperref}
\setlength{\parindent}{0pt}
\setlength{\parskip}{5pt}
\setlength{\emergencystretch}{3em}
\newcommand{\N}{\mathbb{N}}
\newcommand{\Z}{\mathbb{Z}}
\newcommand{\Q}{\mathbb{Q}}
\newcommand{\R}{\mathbb{R}}
\newcommand{\C}{\mathbb{C}}
\newcommand{\F}{\mathbb{F}}
\newcommand{\A}{\mathbb{A}}
\newcommand{\PP}{\mathbb{P}}
\newcommand{\E}{\mathbb{E}}
\newcommand{\eps}{\varepsilon}
\newcommand{\dd}{\,\mathrm{d}}
\begin{document}
\section*{UnsolvedMath ID: 801176}
\textbf{Source problem code:} catalogue entry
\subsection{Definitions and mathematical statement}
For a convex polytope \(\mathcal P\) of dimension \(d\geq 0\), let \(L(\mathcal P)\) be its face lattice, including the empty face \(\hat 0\) and the whole polytope \(\hat 1\). Its rank function satisfies \(\rho(F)=\dim F+1\), so \(\rho(\hat 1)=d+1\). No rationality or simpliciality assumption is imposed.

For the proposed generalization, let \(L\) be a finite graded lattice with least element \(\hat 0\), greatest element \(\hat 1\), and rank \(d+1\). Write \(L^\circ=L\setminus\{\hat 0,\hat 1\}\). The phrase in the source about removing the bottom and top is interpreted as referring to this proper part of a bounded lattice. Let \(\Delta(L^\circ)\) be the simplicial complex whose faces are chains in \(L^\circ\). For a coefficient field \(k\), the Gorenstein\(^*\) condition is interpreted as requiring this complex to be a homology sphere over \(k\), including the corresponding condition on every link: for every face \(\sigma\), including the empty face, and every integer \(j\),
\[
\widetilde H_j\!\left(\operatorname{lk}_{\Delta(L^\circ)}\sigma;k\right)
\cong
\begin{cases}
k,&j=d-1-|\sigma|,\\
0,&j\ne d-1-|\sigma|.
\end{cases}
\]
Here \(|\sigma|\) is the number of vertices of \(\sigma\), and the usual reduced-homology convention for the sphere of dimension \(-1\) is used. The source does not specify the coefficient field or whether it intends a condition over every field; this qualification remains unresolved.

Define the toric, or generalized, \(h\)-polynomial recursively on bounded intervals. For a singleton interval \(Q\), set \(G_Q(t)=1\). For a bounded interval \(Q=[a,b]\) of positive rank \(r\), use its relative rank function \(\rho_Q\) and define
\[
H_Q(t)=\sum_{a\leq x<b}G_{[a,x]}(t)(t-1)^{r-1-\rho_Q(x)}
       =\sum_{i=0}^{r-1}h_i(Q)t^i,
\]
\[
G_Q(t)=\sum_{i=0}^{\lfloor(r-1)/2\rfloor}
\bigl(h_i(Q)-h_{i-1}(Q)\bigr)t^i,
\qquad h_{-1}(Q)=0.
\]
For \(L\) of rank \(d+1\), write \(h(L)=(h_0(L),\ldots,h_d(L))\), using the coefficients of \(H_L(t)\). For a polytope, put \(h(\mathcal P)=h(L(\mathcal P))\).

The source asks whether, for every convex polytope \(\mathcal P\), and more generally for every Gorenstein\(^*\) lattice \(L\) as above, one has
\[
h_0(L)\leq h_1(L)\leq\cdots\leq h_{\lfloor d/2\rfloor}(L),
\]
or, as a weaker requested conclusion,
\[
h_i(L)\geq 0\qquad\text{for every }0\leq i\leq d.
\]
The name Generalized Lower Bound Theorem is used here only for the displayed inequalities requested in the supplied record. The source supplies a dated assessment that the polytope case was solved affirmatively, including nonrational polytopes, while the general Gorenstein\(^*\) lattice case remains open. These are supplied literature-status claims, not independently verified conclusions of this draft.
\subsection{Short English statement}
The question asks whether the toric \(h\)-vector of every convex polytope, or more generally every Gorenstein\(^*\) lattice, increases through its middle entry, or at least has no negative entries. The supplied record labels the problem partially solved and reports an affirmative solution for all convex polytopes, leaving the lattice generalization open. The coefficient convention for Gorenstein\(^*\) remains unspecified.
\subsection{Sources}
\begin{itemize}
\item Supplied local UnsolvedMath registry record: permanent ID \(801176\), source problem code \texttt{catalogue entry}, in the AMR Open Problem Lists collection. Attribution: Stanley, proposed year 1999. Supplied status: \texttt{partially\_solved}. Record updated: 2026-10-16.
\item Stanley, \emph{Positivity Problems and Conjectures in Algebraic Combinatorics} (1999), Problem 2. Supplied source URL: \url{https://math.mit.edu/~rstan/papers/problems.pdf}. The record reports access on 2026-10-16 and extraction by \texttt{pdf+html\_verified}; these are supplied provenance statements, not access or verification performed for this draft.
\item Kalle Karu, \emph{Hard Lefschetz theorem for nonrational polytopes}, Inventiones Mathematicae 157 (2004), 419--447; arXiv identifier \texttt{math/0112087}. The supplied background attributes an affirmative solution for all convex polytopes, rational or nonrational, to this work.
\item Bressler and Lunts, arXiv identifier \texttt{math/0302236}. The supplied background describes this work as a strengthening; it supplies neither a title nor a precise statement of that strengthening.
\item The supplied curation, associated with the record dated 2026-10-16, reports that the general Gorenstein\(^*\) lattice case remains open. This is a supplied dated assessment and has not been independently verified.
\end{itemize}
\end{document}
